\documentclass[11pt,a4paper]{article}

\usepackage{kotex}
\usepackage{fontspec}
\setmainfont{Times New Roman}
\setmainhangulfont{AppleMyungjo}
\setsanshangulfont{Apple SD Gothic Neo}

\usepackage[a4paper,top=18mm,bottom=18mm,left=20mm,right=20mm]{geometry}
\usepackage{array}
\usepackage{enumitem}
\linespread{1.18}
\setlength{\parindent}{0pt}

\newcommand{\sol}[2]{\vspace{6pt}{\sffamily\bfseries #1} \quad 정답 \fbox{#2}\par\vspace{1pt}}

\begin{document}

\begin{center}
{\fontsize{15}{17}\selectfont\sffamily\bfseries 2026 고1 영어 변형 — 지문 29 정답 및 해설}\\[2pt]
{\sffamily 문학의 힘}
\end{center}
\vspace{6pt}

\begin{center}
\renewcommand{\arraystretch}{1.4}
\begin{tabular}{|c|c|c|c|c|}
\hline
\sffamily 문항 & 1 & 2 & 3 & 4 \\\hline
\sffamily 유형 & \sffamily 어법(네모) & \sffamily 빈칸추론 & \sffamily 서술형 & \sffamily 영어 제목 \\\hline
\sffamily 정답 & ① & ③ & (아래) & ① \\\hline
\sffamily 배점 & 2점 & 2점 & 2점 & 2점 \\\hline
\end{tabular}
\end{center}
\vspace{8pt}\hrule\vspace{4pt}

\sol{1. 어법(네모)}{①}
(A) 주어 \textit{It}(Literature)에 이어지는 술어 동사 자리이므로 3인칭 단수 현재형 \textbf{transcends}가 옳다. \textit{transcending}은 준동사로서 단독으로 술어 동사 역할을 할 수 없다.\par
(B) 뒤에 명사구 \textit{the rise of new technologies}가 이어지므로 전치사 \textbf{Despite}(~에도 불구하고)가 옳다. \textit{Although}는 접속사로서 뒤에 주어+동사의 절이 와야 한다.\par
(C) 문학이 현재까지 지속적으로 유지하고 있는 상태를 나타내므로 현재완료 단순형 \textbf{retained}가 자연스럽다. \textit{been retaining}(현재완료 진행형)은 동작의 진행을 강조하여 정적 상태인 \textit{retain}(유지하다)과 어울리지 않는다.\par
따라서 (A) transcends / (B) Despite / (C) retained → 정답 \textbf{①}.

\vspace{4pt}
\sol{2. 빈칸추론}{③}
빈칸이 포함된 문장: \textit{``One reason for literature's enduring power is its ability to capture and convey \underline{\hspace{4cm}}.''}\par
이어지는 문장에서 문학은 ``인간 감정의 깊이(depths of human emotion)'', ``관계의 미묘함(nuances of relationships)'', ``인간 정신의 도전과 승리''를 탐구한다고 설명한다. 이는 모두 인간의 경험이 지닌 복잡성에 해당한다.\par
① 과거 문명의 역사적 사건 — 문학의 핵심 목적과 거리가 멀다.\par
② 창의적 글쓰기에 필요한 기술 — 기술(skill) 측면으로 초점이 어긋난다.\par
\textbf{③ the intricacies of what it means to be human} — ``인간이 된다는 것의 복잡함''으로 the complexities of the human experience의 패러프레이즈. \textbf{정답}.\par
④ 현대 청중의 오락적 필요 — 오락(entertainment)은 본문의 논점이 아니다.\par
⑤ 최신 과학적 발견 — 문학과 무관한 내용이다.

\vspace{4pt}
{\sffamily\bfseries 3. 서술형}\par\vspace{2pt}
\fbox{모범답안} \quad \textbf{having → have}\par\vspace{3pt}
\textbf{채점 기준:} ``have''를 정확히 쓰면 2점 만점. 틀린 단어(\textit{having})의 위치를 찾아 쓴 경우 가산 인정.\par\vspace{3pt}
\textbf{해설:} 주어 \textit{authors}에 이어지는 술어 동사 자리이므로 현재분사 \textit{having} 대신 동사원형(현재형) \textbf{have}가 와야 한다. \textit{having}은 분사로서 단독으로 문장의 서술어가 될 수 없다.

\vspace{8pt}
\sol{4. 제목}{①}
문학은 시대와 공간을 넘어 인간의 생각과 감정을 연결하고, 인간 경험의 복잡성을 포착한다. 따라서 ``인간 경험을 비추는 시대를 초월한 거울''이라는 ①이 가장 적절하다.

\end{document}
